\section{Trader coordinates: the jump-wing handles}
\label{sec:svijw}

The raw quintuple $(a_S,b_S,\rho_S,m_S,s_S)$ is an excellent computational
coordinate system and a poor market language: no single coordinate is
``the level'' or ``the skew,'' and each of $b_S$, $\rho_S$, $s_S$ moves a
wing and the interior at once.  The jump-wing (JW) coordinates of
\citet{GatheralJacquier2014} re-read the same curve through five
functionals chosen to match the desk's questions, and --- the reading this
section presses --- they sort into \emph{three interior handles and two
tail handles}.

\begin{definition}[Jump-wing handles]
\label{def:jw}
At variance time $\tau$, with $w_0=w(0)$ the ATM total variance, the
jump-wing coordinates of a slice are
\begin{equation}
\underbrace{v_J=\frac{w_0}{\tau},\qquad
\psi_J=\big(\sqrt{w}\big)'(0),\qquad
\tilde v_J=\frac{\min_k w(k)}{\tau}}_{\text{interior}},
\qquad
\underbrace{p_J=\frac{\beta_L}{\sqrt{w_0}},\qquad
c_J=\frac{\beta_R}{\sqrt{w_0}}}_{\text{tails}}.
\label{eq:jwdef}
\end{equation}
\end{definition}

Substituting \eqref{eq:svi} gives the closed forms
\begin{equation}
v_J=\frac{w_0}{\tau},\quad
\psi_J=\frac{b_S}{2\sqrt{w_0}}\Big(\rho_S-\frac{m_S}{\sqrt{m_S^2+s_S^2}}\Big),\quad
p_J=\frac{b_S(1-\rho_S)}{\sqrt{w_0}},\quad
c_J=\frac{b_S(1+\rho_S)}{\sqrt{w_0}},\quad
\tilde v_J=\frac{w^*}{\tau},
\label{eq:rawtojw}
\end{equation}
with $w^*$ from \eqref{eq:svivertex}.  Each handle lives in its own units,
and misreading them is the most common JW error.  With
$I(k)=\sqrt{w(k)/\tau}$ the working implied volatility,
\begin{equation}
I(0)=\sqrt{v_J},\qquad
I'(0)=\frac{\psi_J}{\sqrt\tau},\qquad
\min_kI(k)=\sqrt{\tilde v_J},\qquad
\beta_L=p_J\sqrt{v_J\tau},\qquad
\beta_R=c_J\sqrt{v_J\tau}.
\label{eq:jwunits}
\end{equation}
So $v_J$ and $\tilde v_J$ are variances --- the IV chart shows their
square roots; $\psi_J$ is the ATM slope of total \emph{volatility}
$\sqrt w$, and the plotted tangent is $\psi_J/\sqrt\tau$; and $p_J,c_J$
are \emph{normalized} total-variance wing slopes, not slopes of the IV
curve.  \Cref{fig:svihandles} annotates a fitted slice --- SPY December
2026 from the frozen snapshot, the book's recurring node --- in exactly
these units: the fit carries
$\sqrt{v_J}=\MacSviSpyAtmVolPct\%$, tangent slope $\MacSviSpyTangent$,
floor $\sqrt{\tilde v_J}=\MacSviSpyMinVolPct\%$, and wing handles
$(p_J,c_J)=(\MacSviSpyP,\MacSviSpyC)$.

\begin{figure}[htbp]
\centering
\paperfig{0.97\textwidth}{fig_svi_handles.pdf}
\caption{The five handles in their honest units, on the SVI fit of the SPY
December 2026 quotes (rms $\MacSviSpyRmsBp$ vol bp).  (a)~The three
interior handles are read on the IV chart: the level $\sqrt{v_J}$, the
minimum $\sqrt{\tilde v_J}$, and the ATM tangent of slope
$\psi_J/\sqrt\tau$.  (b)~The two tail handles are read on the
total-variance chart, as the slopes $\beta_L=p_J\sqrt{v_J\tau}$ and
$\beta_R=c_J\sqrt{v_J\tau}$ of the asymptotes.  Plotting all five on one
IV axis would give the wing handles units they do not have.}
\label{fig:svihandles}
\end{figure}

One caution before inverting.  The tail handles are \emph{normalized}
indicators: by \eqref{eq:jwunits} the actual asymptotic slopes --- the
objects Lee's bound and the moment map \eqref{eq:momentmap} read --- carry
a factor $\sqrt{v_J\tau}$ of the supposedly interior level.  Raising
$v_J$ with $(p_J,c_J)$ held fixed steepens both actual wings and moves the
inferred moments.  JW is a quoting convention adapted to the desk, not an
orthogonal decomposition of the slice into independent interior and tail
degrees of freedom.

\subsection{Inverting the handles}

Reading \eqref{eq:rawtojw} off a fitted slice is evaluation.  The
interesting direction is the reverse: reconstructing the hyperbola from
five quoted handles.  The two tail handles fix the scale and the tilt,
\begin{equation}
b_S=\frac{\sqrt{w_0}}{2}(p_J+c_J),\qquad
\rho_S=\frac{c_J-p_J}{c_J+p_J},\qquad w_0=v_J\tau ;
\label{eq:invwings}
\end{equation}
the ATM slope fixes the normalized vertex displacement
\begin{equation}
\chi:=\frac{m_S}{\sqrt{m_S^2+s_S^2}}=\rho_S-\frac{4\psi_J}{p_J+c_J},
\label{eq:chidef}
\end{equation}
and the gap between level and floor then sets the width.  Evaluating
\eqref{eq:svi} at $k=0$ and at the vertex gives the two identities
\begin{equation}
w_0=a_S+b_S\Big(\sqrt{m_S^2+s_S^2}-\rho_Sm_S\Big),
\qquad
w^*=a_S+b_S\,s_S\sqrt{1-\rho_S^2};
\label{eq:twoids}
\end{equation}
since $m_S=\chi s_S/\sqrt{1-\chi^2}$ makes
$\sqrt{m_S^2+s_S^2}=s_S/\sqrt{1-\chi^2}$, subtracting the second from the
first leaves $w_0-\tilde v_J\tau=b_S\,s_S\,\mathcal{D}$ with
\begin{equation}
\mathcal{D}(\rho_S,\chi)
=\frac{1-\rho_S\chi}{\sqrt{1-\chi^2}}-\sqrt{1-\rho_S^2},
\label{eq:Ddef}
\end{equation}
so that
\begin{equation}
s_S=\frac{(v_J-\tilde v_J)\,\tau}{b_S\,\mathcal{D}},\qquad
m_S=\frac{\chi\,s_S}{\sqrt{1-\chi^2}},\qquad
a_S=\tilde v_J\tau-b_S\,s_S\sqrt{1-\rho_S^2}.
\label{eq:jwinverse}
\end{equation}
The denominator's sign is Cauchy--Schwarz.  Applied to the unit vectors
$(\rho_S,\sqrt{1-\rho_S^2})$ and $(\chi,\sqrt{1-\chi^2})$,
\begin{equation}
\rho_S\chi+\sqrt{1-\rho_S^2}\sqrt{1-\chi^2}\;\le\;1,
\quad\text{i.e.}\quad
\mathcal{D}\;\ge\;0,
\label{eq:Dpos}
\end{equation}
with equality exactly when the vectors coincide, $\chi=\rho_S$.  This
single inequality gives both the domain and the singularity of the
coordinate system.

\begin{theorem}[The image of the jump-wing map]
\label{thm:jwimage}
Fix $\tau>0$ and restrict \eqref{eq:svi} to $b_S>0$, $|\rho_S|<1$,
$s_S>0$, $w_0>0$.
\begin{enumerate}[label=(\roman*)]
\item A JW point has a unique inverse, given by
\eqref{eq:invwings}--\eqref{eq:jwinverse}, exactly on
\begin{equation}
v_J>0,\qquad p_J>0,\ c_J>0,\qquad
-\frac{p_J}{2}<\psi_J<\frac{c_J}{2},\qquad
\psi_J\ne0,\qquad \tilde v_J<v_J .
\label{eq:regdom}
\end{equation}
\item The image also contains the stratum $\psi_J=0$, $\tilde v_J=v_J$
(with $v_J,p_J,c_J>0$), every point of which is attained by infinitely
many slices.  There are no other image points.
\end{enumerate}
\end{theorem}

\begin{proof}
If $p_J,c_J>0$ then \eqref{eq:invwings} gives $b_S>0$ and
$\rho_S\in(-1,1)$; conversely every admissible slice has $p_J,c_J>0$.
From \eqref{eq:chidef}, using $(\rho_S-1)(p_J+c_J)=-2p_J$ and
$(\rho_S+1)(p_J+c_J)=2c_J$,
\[
|\chi|<1
\iff \rho_S-1<\frac{4\psi_J}{p_J+c_J}<\rho_S+1
\iff -\frac{p_J}{2}<\psi_J<\frac{c_J}{2},
\]
which keeps $m_S$ finite.  By \eqref{eq:Dpos}, $\mathcal{D}>0$ exactly
when $\chi\ne\rho_S$, i.e.\ $\psi_J\ne0$; then $s_S>0$ in
\eqref{eq:jwinverse} requires $v_J>\tilde v_J$.  Each failure breaks a
requirement: one of $p_J,c_J\le0$ forces $b_S\le0$ or $|\rho_S|\ge1$;
$\psi_J$ outside the band sends $|\chi|\ge1$; $\psi_J=0$ annihilates
$\mathcal{D}$; $\tilde v_J\ge v_J$ makes $s_S\le0$.  For the stratum:
$\psi_J=0$ gives $\chi=\rho_S$, which by \eqref{eq:svivertex} puts the
vertex at $k^*=0$, so the floor is attained at the money and
$\tilde v_J=v_J$ automatically.  Conversely fix $(v_J,p_J,c_J)$ on the
stratum; \eqref{eq:invwings} sets $(b_S,\rho_S)$, and for \emph{every}
$s_S>0$ the slice with $m_S=\rho_S s_S/\sqrt{1-\rho_S^2}$ and
$a_S=v_J\tau-b_Ss_S\sqrt{1-\rho_S^2}$ reproduces the same five handles.
\end{proof}

\begin{remark}[The blind spot is the interior curvature]
\label{rem:stratum}
The free width on the stratum is not an abstract degeneracy.  There the
vertex sits at the money, the handles pin the level and both tails, and
what they omit is exactly how sharply the curve turns through its bottom:
by \eqref{eq:svivertex} the vertex curvature $\kappa^*=b_Sq_S^3/s_S$ runs
through every positive value as $s_S$ does.  \Cref{fig:svistratum}(b)
draws three slices sharing one handle quintuple.  Nor is the failure
confined to the stratum itself: $\mathcal{D}$ vanishes at $\chi=\rho_S$
together with its first $\chi$-derivative, and a direct differentiation
of \eqref{eq:Ddef} gives
$\partial^2\mathcal{D}/\partial\chi^2=(1-\rho_S^2)^{-3/2}$ there, so to
second order,
\begin{equation}
\mathcal{D}
\sim\frac{(\chi-\rho_S)^2}{2\,(1-\rho_S^2)^{3/2}}
=\frac{8\,\psi_J^2}{(p_J+c_J)^2(1-\rho_S^2)^{3/2}},
\qquad\psi_J\to0,
\label{eq:Dasym}
\end{equation}
so the denominator vanishes \emph{quadratically}: near the stratum,
recovering a finite width forces $v_J-\tilde v_J=O(\psi_J^2)$, quote noise
in the handles is amplified by $1/\psi_J^2$, and a direct evaluation of
\eqref{eq:Ddef} subtracts two nearly equal numbers.  An implementation can
remove the cancellation by evaluating an algebraically identical form; no
rearrangement can restore a curvature the handles never carried.  The
stratum is not exotic --- ATM sitting at the smile's minimum is where a
desk would expect the coordinates to be simplest, and fitted index slices
pass close to it.
\end{remark}

\begin{figure}[htbp]
\centering
\paperfig{0.97\textwidth}{fig_svi_stratum.pdf}
\caption{Where the jump-wing chart tears.  (a)~The inverse denominator
$\mathcal{D}$ against $\psi_J$ at fixed $(p_J,c_J)$: it closes
quadratically at $\psi_J=0$, following \eqref{eq:Dasym} (dashed).
(b)~Three slices with identical handles $(v_J,0,p_J,c_J,v_J)$ --- same
level, same asymptotes --- differing only in width $s_S$: the five
handles lose the interior curvature (handle agreement across the three
slices: $\MacSviStratumSpread$).}
\label{fig:svistratum}
\end{figure}

The screens of \cref{sec:svibelly} translate directly into this
vocabulary, which is convenient for checking a quoted handle set before
any conversion: the floor condition is $\tilde v_J\ge0$; the domain
condition is $-p_J/2<\psi_J<c_J/2$; and since
$b_S(1+|\rho_S|)=\max(\beta_L,\beta_R)=\max(p_J,c_J)\sqrt{w_0}$ by
\eqref{eq:invwings} and \eqref{eq:jwunits}, the cap \eqref{eq:leecap}
reads
\begin{equation}
\max(p_J,c_J)\,\sqrt{v_J\tau}\;\le\;\beta_{\max}.
\label{eq:jwcap}
\end{equation}
None of the three certifies the interior, for the reason
\cref{sec:svibelly} established.
